

\chapter{Tools}
\label{sec:Tools}

The tools described in the following are intended to facilitate the handling of \textit{ms}2. They should allow for an easy start of molecular simulations with \textit{ms}2 and give access to all output data generated by \textit{ms}2.


\section{\textit{ms}2par}
\label{sec:ms2par}


The GUI-based feature tool $ms2par$ allows for a convenient generation of *.par files according to user specifications, cf. Figure~\ref{pic_ms2par}. It assigns the molecular model, ensemble, thermodynamic state point, time step, number of equilibration steps, etc. Furthermore, the tool proposes a cut-off radius based on the input data. After the user completes the specifications, the program generates the file *.par in ASCII format, to be read by $ms2$. Applying $ms2par$ is simple and should be intuitive even for new users. $ms2par$ is a java application, thus it runs independently on all operating systems. 


\section{\textit{ms}2potmod}
\label{sec:ms2potmod}



\section{\textit{ms}2chart}
\label{sec:ms2chart}


$ms2chart$ is a GUI-based java applet for evaluating the simulation results from $ms2$, cf. Figure~\ref{pic_ms2chart}. This tool plots the evolution of properties calculated by $ms2$. The properties shown on the different graph axes are chosen from a drop 
down menu and the plot is shown directly in the GUI. \\
For a better analysis, $ms2chart$ allows various features: plotting block averages as well as simulation averages into the same plot, changing the design of the plot, individual labeling of the axes and adjusting the frame detail of the plot. All plots can be exported in png~format. If needed, the simulation details can be displayed, i.e. the specifications in the *.par file, the results in the~*.res file and the summary in the~*.log file. \\
The analysis program $ms2chart$ can be executed at any time of the simulation, i.e. also while the simulation is still ongoing. There is no loss of data, if it is run on the fly. The tool $ms2chart$ is easy to understand and can be handled intuitively. It allows for an easy first analysis of the $ms2$ simulation results. Advanced users will use $ms2chart$ for a quick check of their simulation results, such as whether the equilibration time was appropriate or for a fast analysis, which is important if extensive series of simulation runs are performed.

\pagebreak

\section{\textit{ms}2molecules}
\label{sec:ms2molecules}
 

The program $ms2molecules$ is a visualization tool for $ms2$. The program displays the molecular trajectories that are stored 
by $ms2$ as a series of configurations in the *.vim file. $ms2molecules$ visualizes molecular sites by colored spheres. The colors are user defined. The size of a sphere is by default proportional to the LJ~size parameter $\sigma$, but can be changed manually in the first lines of the *.vim-file. This feature facilitates monitoring a component of interest by reducing the size of other components, e.g. solvents. Other features like zooming in or out of the simulation volume and rotating the simulation box also facilitates the analysis of the simulation. The visualization can be exported via snapshots in the jpg~format. \\
The program is based on OpenGL and written in C. The handling of the program is simple and console based. Requirements for the feature tool are OpenGL in a Windows or Linux environment. Figure~\ref{pic_ms2molecules} shows a snapshot of a ternary mixture, taken with the program $ms2molecules$, which is convenient to gain insight into the trajectory of the system, or the state of the system, respectively.\\

\textbf{Colors:}\\
Setting the colours of the molecules in the visualisation window has to be done manually in the header of the \texttt{*.vim}-file (see section \ref{sec:vim - file}). The particle colour is defined by the last integer value of each line in the header, which makes the visualisation easy ajustable. The colour of each cite-type of each molecule-type in the simulation can be set separately. The available colours are:
\begin{table}[htbp]
	\centering
	\caption{Available Colours in \textit{molecules}}
	\begin{tabular}{rl}
		\toprule
		Integer in vim file & Colour \\
		\midrule
		1     & black \\
		2     & red \\
		3     & orange \\
		4     & yellow \\
		5     & yellow 2 \\
		6     & green \\
		7     & green 2 \\
		8     & turquoise \\
		9     & blue \\
		10    & blue 2 \\
		11    & purple \\
		12    & pink \\
		13    & pink 2 \\
		14    & grey \\
		\bottomrule
	\end{tabular}%
	\label{tab:Colours in vim file}%
\end{table}%




\textbf{Adjusting Atom size:}\\
The atom size, or the size of one cite in a molecule-type can be adjusted manually by the user in the header of the \texttt{*vim}-file. That can be done by setting the sigma-values of each cite, e.\,g.
\begin{Verbatim}
~  1 LJ  0.0000 -0.3387 -1.4713  3.0  1
~  1 LJ -0.0000  0.8836  0.0917  1.42  2
~  1 LJ -0.0000 -0.4503  1.1711  1.8  3 
\end{Verbatim}
The integer "1" in the first column after the tilde, specifies the LJ-cites as one of the 1st molecule (used \texttt{*pm} file). The 2nd, 3rd and 4th numeric column are the coordinates of the cite relative to the centre of mass of the molecule. The 5th column indicates the sigma value of the cite and can manually re-adjusted. The last integer column indicates the colour of the cite (see above). 

\pagebreak

\section{res2gepar}
\label{sec:res2gepar}
% Minh







